\section{共享年表之一：从三家分晋到各国求变}

\begin{event}{公元前453年}{韩赵魏灭智氏}{可靠纪年}{晋}\eventanchor{WQ-0453-ZHI}{战国·三家灭智}
\term{这一年发生了什么}：韩、赵、魏联合消灭智氏，晋国公室已无法控制主要卿族。
\term{后来改变了什么}：战国的一个起点，不是新国家突然出现，而是春秋大国晋的内部资源被几个家族分割。
\end{event}

\begin{event}{公元前403年}{韩赵魏获周王室册命}{可靠纪年}{洛邑}\eventanchor{WQ-0403-THREEJIN}{战国·三家分晋}
\term{这一年发生了什么}：周威烈王承认韩、赵、魏为诸侯。
\term{后来改变了什么}：周王室承认既成事实，说明名分仍有价值，但已无法逆转实力造成的政治重组。
\end{event}

\begin{event}{公元前356年}{商鞅第一次变法}{可靠纪年}{秦}\eventanchor{WQ-0356-SHANGYANG}{战国·商鞅变法}
\term{这一年发生了什么}：秦孝公任用商鞅，传统叙事将什伍连坐、军功爵、编户和基层控制等改革与此相连。
\term{后来改变了什么}：秦把人口、土地、军功与行政更紧地连在一起；但它不是战国唯一的改革者，必须放回魏、楚、韩等国的竞争中理解。
\end{event}
